% GENERATED FILE. Edit canonical lessons, then run npm run generate:books.
% canonical-source-hash: fnv1a64:0a26350bbd94059a
% canonical-lessons: RU-C435-obsluzhivaniye, RU-C435-obsluzhivat, RU-C435-personal, RU-C435-chayevyye, RU-C435-khamit

\chapter{Service and tips}
\label{ch:ru-service-and-tips}
\hlchaptermodality{435}

\begin{chapteropening}
\textbf{By the end of this chapter:} \emph{I can talk about service: the service, serving customers, the staff, tips and rudeness.}

Complete the last lesson of chapter 435: I can talk about service: the service, serving customers, the staff, tips and rudeness.
\end{chapteropening}

\section[\texorpdfstring{\ru{обслуживание} (obslúzhivaniye)}{обслуживание (obslúzhivaniye)}]{\ru{обслуживание} (obslúzhivaniye) — service}

\label{lesson:RU-C435-obsluzhivaniye}



\begin{quote}
Before the new one: say the Russian for dissatisfied, then the Russian for out of order.
\end{quote}

\subsection*{You'll want to know: \ru{обслуживание}}
\textbf{\ru{обслуживание}} (\emph{obslúzhivaniye}) — \textquotedblleft{}service\textquotedblright{}.

\begin{quote}
\ru{В} \ru{этой} \ru{гостинице} \ru{прекрасное} \ru{обслуживание}.
\end{quote}

\begin{quote}
\emph{V étoy gostínitse prekrásnoye obslúzhivaniye.}
\end{quote}

\begin{quote}
This hotel has excellent service.
\end{quote}

Neuter. \emph{\ru{Обслуживание} \ru{номеров}} — room service.

\begin{cousinweb}
\emph{\ru{обслуживание}} is built on \emph{\ru{обслуживать}} 'to serve', from \emph{\ru{об}-} and \emph{\ru{служить}} 'to serve'; EtymDB lists both without a deeper source.
\end{cousinweb}

\begin{grammarlens}[title={What you've built}]
One more word for telling, explaining and describing at B1.
\end{grammarlens}

\subsection*{Your turn}
\begin{itemize}
\raggedright
  \item \emph{Say it:} \emph{obslúzhivaniye}
  \item \emph{Say it:} \emph{obslúzhivaniye}, once more
  \item \emph{Say it:} a sentence of your own with \emph{obslúzhivaniye}
  \item \emph{Recall:} say the Russian for dissatisfied, then the Russian for out of order, then say \emph{obslúzhivaniye} again
\end{itemize}

\begin{tcolorbox}[breakable,colback=teal!4,colframe=teal!35!black,title={Before you move on}]
What does \ru{обслуживание} mean? (\textquotedblleft{}Service\textquotedblright{}.) Say it once more.
\end{tcolorbox}
\section[\texorpdfstring{\ru{обслуживать} (obslúzhivat)}{обслуживать (obslúzhivat)}]{\ru{обслуживать} (obslúzhivat) — to serve (customers); to service (a car)}

\label{lesson:RU-C435-obsluzhivat}



\begin{quote}
Before the new one: say the Russian for out of order, then the Russian for service.
\end{quote}

\subsection*{You'll want to know: \ru{обслуживать}}
\textbf{\ru{обслуживать}} (\emph{obslúzhivat}) — \textquotedblleft{}to serve (customers); to service (a car)\textquotedblright{}.

\begin{quote}
\ru{Официант} \ru{обслуживает} \ru{пять} \ru{столов}.
\end{quote}

\begin{quote}
\emph{Ofitsiánt obslúzhivayet pyat stolóv.}
\end{quote}

\begin{quote}
The waiter serves five tables.
\end{quote}

Imperfective; perfective \emph{\ru{обслужить}}.

\begin{cousinweb}
\emph{\ru{обслуживать}} is built from \emph{\ru{об}-} and \emph{-\ru{служивать}}, from \emph{\ru{служить}} 'to serve'; EtymDB lists \emph{\ru{служить}} without a deeper source.
\end{cousinweb}

\begin{grammarlens}[title={What you've built}]
One more word for telling, explaining and describing at B1.
\end{grammarlens}

\subsection*{Your turn}
\begin{itemize}
\raggedright
  \item \emph{Say it:} \emph{obslúzhivat}
  \item \emph{Say it:} \emph{obslúzhivat}, once more
  \item \emph{Say it:} a sentence of your own with \emph{obslúzhivat}
  \item \emph{Recall:} say the Russian for out of order, then the Russian for service, then say \emph{obslúzhivat} again
\end{itemize}

\begin{tcolorbox}[breakable,colback=teal!4,colframe=teal!35!black,title={Before you move on}]
What does \ru{обслуживать} mean? (\textquotedblleft{}To serve (customers); to service (a car)\textquotedblright{}.) Say it once more.
\end{tcolorbox}
\section[\texorpdfstring{\ru{персонал} (personál)}{персонал (personál)}]{\ru{персонал} (personál) — staff, personnel}

\label{lesson:RU-C435-personal}



\begin{quote}
Before the new one: say the Russian for service, then the Russian for to serve.
\end{quote}

\subsection*{You'll want to know: \ru{персонал}}
\textbf{\ru{персонал}} (\emph{personál}) — \textquotedblleft{}staff, personnel\textquotedblright{}.

\begin{quote}
\ru{Персонал} \ru{гостиницы} \ru{был} \ru{очень} \ru{вежливый}.
\end{quote}

\begin{quote}
\emph{Personál gostínitsy byl óchen vézhlivyy.}
\end{quote}

\begin{quote}
The hotel staff were very polite.
\end{quote}

Masculine and singular in Russian: \emph{\ru{персонал} \ru{был}}.

\begin{cousinweb}
\emph{\ru{персонал}} is usually traced to German \emph{Personal} 'staff', from Medieval Latin \emph{persōnālis} 'of a person', on \emph{persōna} 'a mask, a person' (Lewis \& Short). EtymDB does not list it.
\end{cousinweb}

\begin{grammarlens}[title={What you've built}]
One more word for telling, explaining and describing at B1.
\end{grammarlens}

\subsection*{Your turn}
\begin{itemize}
\raggedright
  \item \emph{Say it:} \emph{personál}
  \item \emph{Say it:} \emph{personál}, once more
  \item \emph{Say it:} a sentence of your own with \emph{personál}
  \item \emph{Recall:} say the Russian for service, then the Russian for to serve, then say \emph{personál} again
\end{itemize}

\begin{tcolorbox}[breakable,colback=teal!4,colframe=teal!35!black,title={Before you move on}]
What does \ru{персонал} mean? (\textquotedblleft{}Staff, personnel\textquotedblright{}.) Say it once more.
\end{tcolorbox}
\section[\texorpdfstring{\ru{чаевые} (chayevýye)}{чаевые (chayevýye)}]{\ru{чаевые} (chayevýye) — a tip, tips (for service)}

\label{lesson:RU-C435-chayevyye}



\begin{quote}
Before the new one: say the Russian for to serve, then the Russian for staff.
\end{quote}

\subsection*{You'll want to know: \ru{чаевые}}
\textbf{\ru{чаевые}} (\emph{chayevýye}) — \textquotedblleft{}a tip, tips (for service)\textquotedblright{}.

\begin{quote}
\ru{В} \ru{ресторане} \ru{мы} \ru{оставили} \ru{чаевые}.
\end{quote}

\begin{quote}
\emph{V restoráne my ostávili chayevýye.}
\end{quote}

\begin{quote}
We left a tip in the restaurant.
\end{quote}

Plural only; an adjective used as a noun: money \emph{\ru{на} \ru{чай}}, for tea.

\begin{cousinweb}
\emph{\ru{чаевые}} is built on \emph{\ru{чай}} 'tea' (EtymDB lists \emph{\ru{чай}}): money once given \emph{\ru{на} \ru{чай}}, for a glass of tea.
\end{cousinweb}

\begin{grammarlens}[title={What you've built}]
One more word for telling, explaining and describing at B1.
\end{grammarlens}

\subsection*{Your turn}
\begin{itemize}
\raggedright
  \item \emph{Say it:} \emph{chayevýye}
  \item \emph{Say it:} \emph{chayevýye}, once more
  \item \emph{Say it:} a sentence of your own with \emph{chayevýye}
  \item \emph{Recall:} say the Russian for to serve, then the Russian for staff, then say \emph{chayevýye} again
\end{itemize}

\begin{tcolorbox}[breakable,colback=teal!4,colframe=teal!35!black,title={Before you move on}]
What does \ru{чаевые} mean? (\textquotedblleft{}A tip, tips (for service)\textquotedblright{}.) Say it once more.
\end{tcolorbox}
\section[\texorpdfstring{\ru{хамить} (khamít)}{хамить (khamít)}]{\ru{хамить} (khamít) — to be rude (to someone)}

\label{lesson:RU-C435-khamit}



\begin{quote}
Before the new one: say the Russian for staff, then the Russian for a tip.
\end{quote}

\subsection*{You'll want to know: \ru{хамить}}
\textbf{\ru{хамить}} (\emph{khamít}) — \textquotedblleft{}to be rude (to someone)\textquotedblright{}.

\begin{quote}
\ru{Официант} \ru{хамил} \ru{всем} \ru{посетителям}.
\end{quote}

\begin{quote}
\emph{Ofitsiánt khamíl vsem posetítelyam.}
\end{quote}

\begin{quote}
The waiter was rude to all the customers.
\end{quote}

Imperfective; perfective \emph{\ru{нахамить}}. Takes the dative. Present \emph{\ru{хамлю}, \ru{хамишь}}. Colloquial.

\begin{cousinweb}
\emph{\ru{хамить}} is built on \emph{\ru{хам}} 'a boor'. \emph{\ru{Хам}} is usually traced to a name from the Bible, a son who mocked his father; EtymDB lists it without a source, so the line is hedged.
\end{cousinweb}

\begin{grammarlens}[title={What you've built}]
One more word for telling, explaining and describing at B1.
\end{grammarlens}

\subsection*{Your turn}
\begin{itemize}
\raggedright
  \item \emph{Say it:} \emph{khamít}
  \item \emph{Say it:} \emph{khamít}, once more
  \item \emph{Say it:} a sentence of your own with \emph{khamít}
  \item \emph{Recall:} say the Russian for staff, then the Russian for a tip, then say \emph{khamít} again
\end{itemize}

\begin{tcolorbox}[breakable,colback=teal!4,colframe=teal!35!black,title={Before you move on}]
What does \ru{хамить} mean? (\textquotedblleft{}To be rude (to someone)\textquotedblright{}.) Say it once more.
\end{tcolorbox}
